\documentclass[11pt,a4paper]{article}
\usepackage[utf8]{inputenc}
\usepackage[spanish]{babel}
\usepackage{amsmath,amssymb,amsfonts}
\usepackage{geometry}
\usepackage{xcolor}
\usepackage{enumitem}
\usepackage{tcolorbox}
\usepackage{hyperref}

\geometry{top=2.5cm,left=2.5cm,right=2.5cm,bottom=2.5cm}
\setlist{leftmargin=*}
\pagestyle{plain}
\setlength{\parindent}{0pt}
\setlength{\parskip}{6pt}

\definecolor{azuloscuro}{RGB}{0,51,102}
\definecolor{azulclaro}{RGB}{51,153,255}
\definecolor{rojoclaro}{RGB}{255,102,102}
\definecolor{verdeclaro}{RGB}{102,204,102}
\definecolor{amarillo}{RGB}{255,230,153}

\hypersetup{
    colorlinks=true,
    linkcolor=azuloscuro,
    urlcolor=azuloscuro,
    citecolor=azuloscuro,
    pdfborder={0 0 0}
}

\title{\textbf{\LARGE GUÍA COMPLETA DE COMUNICACIONES DIGITALES}\\[0.5cm]
\Large Teoría Explicada con Detalle de Cada Fórmula\\[0.3cm]
\large TAF STAR UEC B - IMT Atlantique 2024-2025}
\author{}
\date{}

\begin{document}

\maketitle
\thispagestyle{empty}

\vspace{1cm}

\begin{center}
\fbox{\begin{minipage}{0.8\textwidth}
\textbf{Objetivo de esta guía:}

Este documento explica detalladamente cada fórmula y concepto de las comunicaciones digitales. Para cada ecuación encontrarás:
\begin{itemize}
\item Significado de cada variable y símbolo
\item Interpretación física del concepto
\item Por qué es importante
\item Aplicaciones prácticas
\end{itemize}
\end{minipage}}
\end{center}

\newpage

\tableofcontents

\newpage

\section{Fundamentos de las Comunicaciones Digitales}

\subsection{¿Qué son las comunicaciones digitales?}

Las comunicaciones digitales permiten transmitir información discreta (bits: 0 y 1) a través de un medio físico continuo (ondas electromagnéticas, señales eléctricas, luz en fibra óptica).

\textbf{Problema fundamental:} Necesitamos enviar datos digitales usando señales analógicas, minimizando errores y maximizando la velocidad de transmisión.

\subsection{Cadena de comunicación completa}

\textbf{Emisor:}
\begin{enumerate}
\item \textbf{Codificación de fuente:} Compresión (eliminar redundancia innecesaria)
\item \textbf{Codificación de canal:} Añadir redundancia controlada para detectar/corregir errores
\item \textbf{Modulación:} Convertir bits en señales analógicas que se pueden transmitir
\item \textbf{Filtrado:} Adaptar la señal al canal (limitar ancho de banda)
\end{enumerate}

\textbf{Canal físico:}
\begin{itemize}
\item Añade \textbf{ruido aleatorio} (AWGN - Additive White Gaussian Noise)
\item Introduce \textbf{interferencias} de otras señales
\item Causa \textbf{distorsión} (multipaso, eco, reflexiones)
\item Produce \textbf{atenuación} (pérdida de potencia)
\end{itemize}

\textbf{Receptor:}
\begin{enumerate}
\item \textbf{Filtro adaptado:} Maximiza la relación señal/ruido (SNR)
\item \textbf{Muestreo:} Convierte señal continua en muestras discretas
\item \textbf{Decisión:} Determina qué símbolo se transmitió (estima bits)
\item \textbf{Decodificación:} Detecta y corrige errores usando redundancia
\end{enumerate}

\textbf{Objetivo global:} Transmitir la mayor cantidad de bits por segundo (débito) con la menor tasa de error posible, usando el mínimo ancho de banda y potencia.

\newpage

\section{Transformada de Fourier - Herramienta Fundamental}

\subsection{Concepto y motivación}

La Transformada de Fourier (FT) es esencial porque permite analizar señales en dos dominios complementarios:

\begin{itemize}
\item \textbf{Dominio temporal $h(t)$:} Cómo varía la señal en el tiempo
\item \textbf{Dominio frecuencial $H(\nu)$:} Qué frecuencias contiene la señal y con qué amplitud
\end{itemize}

\textbf{¿Por qué es importante en comunicaciones?}
\begin{itemize}
\item El ancho de banda del canal limita qué frecuencias pueden pasar
\item El filtrado es multiplicación en frecuencia (más simple que convolución en tiempo)
\item Nos permite diseñar señales que usen eficientemente el espectro disponible
\item Podemos analizar el efecto del ruido y la distorsión en cada frecuencia
\end{itemize}

\subsection{Definición de la Transformada de Fourier}

\begin{tcolorbox}[colback=verdeclaro!20,colframe=azuloscuro,title=\textbf{Transformada de Fourier Directa}]
$$\boxed{H(\nu) = \int_{\mathbb{R}} h(t)e^{-i2\pi\nu t}dt}$$

\textbf{Cada símbolo significa:}
\begin{itemize}
\item $h(t)$: señal en el dominio del tiempo (función de $t$)
\item $H(\nu)$: señal en el dominio de la frecuencia (función de $\nu$)
\item $t$: tiempo [segundos]
\item $\nu$: frecuencia [Hz = ciclos/segundo]
\item $i$: unidad imaginaria ($i^2 = -1$)
\item $e^{-i2\pi\nu t} = \cos(2\pi\nu t) - i\sin(2\pi\nu t)$: exponencial compleja
\item $\int_{\mathbb{R}} \ldots dt$: integral sobre todo el eje temporal
\end{itemize}

\textbf{¿Qué hace esta fórmula?}

Descompone la señal temporal $h(t)$ en una suma infinita de oscilaciones sinusoidales de diferentes frecuencias. El valor $H(\nu)$ (que es un número complejo) nos dice:
\begin{itemize}
\item \textbf{Módulo $|H(\nu)|$:} Cuánta "cantidad" de frecuencia $\nu$ hay en la señal
\item \textbf{Fase $\arg(H(\nu))$:} El desfase de esa componente frecuencial
\end{itemize}

\textbf{Interpretación física:}

La exponencial $e^{-i2\pi\nu t}$ "interroga" la señal $h(t)$ buscando componentes que oscilen a frecuencia $\nu$. Si hay mucha energía a esa frecuencia, $H(\nu)$ será grande. Si no hay nada a esa frecuencia, $H(\nu) \approx 0$.
\end{tcolorbox}

\begin{tcolorbox}[colback=verdeclaro!20,colframe=azuloscuro,title=\textbf{Transformada Inversa de Fourier}]
$$\boxed{h(t) = \int_{\mathbb{R}} H(\nu)e^{i2\pi\nu t}d\nu}$$

\textbf{¿Qué hace?}

Reconstruye la señal temporal $h(t)$ sumando (integrando) todas las componentes frecuenciales $H(\nu)$ multiplicadas por sus respectivas oscilaciones $e^{i2\pi\nu t}$.

\textbf{Interpretación fundamental:}

Cualquier señal $h(t)$, por compleja que sea, puede representarse como una suma de exponenciales complejas (senos y cosenos) con amplitudes y fases específicas dadas por $H(\nu)$.

\textbf{Ejemplo concreto:}

Si $h(t) = \cos(2\pi f_0 t)$, entonces $H(\nu)$ tendrá dos "picos" (deltas de Dirac) en $\nu = \pm f_0$. Es decir, solo hay energía en esas dos frecuencias.
\end{tcolorbox}

\subsection{Propiedades esenciales de la Transformada de Fourier}

\subsubsection{Propiedad 1: Convolución en tiempo = Producto en frecuencia}

\begin{tcolorbox}[colback=rojoclaro!20,colframe=azuloscuro,title=\textbf{Convolución $\leftrightarrow$ Producto}]
$$\boxed{\text{FT}[h(t) * g(t)] = H(\nu) \cdot G(\nu)}$$

\textbf{Definición de convolución:}
$$h(t) * g(t) = \int_{-\infty}^{\infty} h(\tau)g(t-\tau)d\tau$$

\textbf{Significado de cada término:}
\begin{itemize}
\item $h(t) * g(t)$: convolución de $h$ y $g$ (operación "pesada" en tiempo)
\item $\tau$: variable de integración (tiempo auxiliar)
\item $H(\nu)G(\nu)$: simple multiplicación punto a punto en frecuencia
\end{itemize}

\textbf{¿Por qué es tan importante?}

El filtrado en tiempo es una convolución: si una señal $h(t)$ pasa por un filtro con respuesta impulsional $g(t)$, la salida es $y(t) = h(t) * g(t)$.

En frecuencia, esto se convierte en una simple multiplicación: $Y(\nu) = H(\nu)G(\nu)$. Es decir:
\begin{itemize}
\item En tiempo: convolución (operación compleja, muchas multiplicaciones y sumas)
\item En frecuencia: producto (operación trivial, una multiplicación por cada $\nu$)
\end{itemize}

\textbf{Aplicación práctica:}

Cuando diseñamos filtros, es mucho más fácil trabajar en frecuencia: simplemente multiplicamos el espectro de entrada por la función de transferencia del filtro.
\end{tcolorbox}

\subsubsection{Propiedad 2: Producto en tiempo = Convolución en frecuencia}

\begin{tcolorbox}[colback=rojoclaro!20,colframe=azuloscuro,title=\textbf{Producto $\leftrightarrow$ Convolución}]
$$\boxed{\text{FT}[h(t) \cdot g(t)] = H(\nu) * G(\nu)}$$

\textbf{Significado:}

Multiplicar dos señales en tiempo → convolucionar sus espectros en frecuencia.

\textbf{¿Por qué es importante?}

Esta propiedad explica lo que ocurre en la modulación. Cuando multiplicamos una señal banda base por una portadora, estamos multiplicando en tiempo, lo que causa una convolución (mezcla) en frecuencia.

\textbf{Ejemplo de modulación AM:}

Si multiplicamos $h(t)$ por $\cos(2\pi f_c t)$ (portadora), el espectro $H(\nu)$ se replica alrededor de $\pm f_c$. La señal que estaba en banda base (frecuencias bajas) ahora está centrada en la frecuencia de la portadora $f_c$ (frecuencias altas).
\end{tcolorbox}

\subsubsection{Propiedad 3: Desplazamiento temporal}

\begin{tcolorbox}[colback=rojoclaro!20,colframe=azuloscuro,title=\textbf{Retardo en tiempo}]
$$\boxed{\text{FT}[h(t-a)] = H(\nu)e^{-i2\pi\nu a}}$$

\textbf{Variables:}
\begin{itemize}
\item $a$: retardo temporal [segundos] (constante)
\item $h(t-a)$: señal retrasada $a$ segundos
\item $e^{-i2\pi\nu a}$: factor de fase que depende linealmente de $\nu$
\end{itemize}

\textbf{Significado físico:}

Un retardo en tiempo solo afecta la \textbf{fase} del espectro, no su \textbf{magnitud}.

$$|H(\nu)e^{-i2\pi\nu a}| = |H(\nu)| \cdot |e^{-i2\pi\nu a}| = |H(\nu)| \cdot 1 = |H(\nu)|$$

La magnitud del espectro permanece igual.

\textbf{Consecuencia importante:}

La energía espectral (proporcional a $|H(\nu)|^2$) no cambia con retardos. Solo cambia "cuándo" llega la señal, no "qué frecuencias" contiene.

\textbf{Aplicación:}

En canales multitrayecto, diferentes copias de la señal llegan con diferentes retardos. Cada copia tiene el mismo espectro de magnitud pero diferente fase → interferencia constructiva o destructiva dependiendo de la frecuencia.
\end{tcolorbox}

\subsubsection{Propiedad 4: Modulación (desplazamiento frecuencial)}

\begin{tcolorbox}[colback=amarillo!40,colframe=azuloscuro,title=\textbf{BASE MATEMÁTICA DE LA MODULACIÓN}]
$$\boxed{\text{FT}[h(t)e^{i2\pi\nu_0 t}] = H(\nu-\nu_0)}$$

\textbf{Variables y términos:}
\begin{itemize}
\item $\nu_0$: frecuencia de la portadora [Hz]
\item $e^{i2\pi\nu_0 t} = \cos(2\pi\nu_0 t) + i\sin(2\pi\nu_0 t)$: oscilación compleja a frecuencia $\nu_0$
\item $h(t)e^{i2\pi\nu_0 t}$: señal modulada (banda base multiplicada por portadora)
\item $H(\nu-\nu_0)$: espectro original desplazado $\nu_0$ Hz a la derecha
\end{itemize}

\textbf{Significado físico profundo:}

Multiplicar una señal por una exponencial compleja en tiempo \textbf{traslada} su espectro en frecuencia.

\textbf{¿Qué pasa visualmente?}

Si $H(\nu)$ estaba centrado en $\nu = 0$ (banda base), después de multiplicar por $e^{i2\pi\nu_0 t}$, el espectro se mueve íntegramente a estar centrado en $\nu = \nu_0$.

\textbf{Aplicación fundamental - La modulación:}

Para transmitir una señal banda base $h(t)$ (que ocupa frecuencias bajas, por ejemplo 0-5 kHz) en una frecuencia portadora alta $\nu_0$ (por ejemplo 100 MHz en FM), la multiplicamos por $e^{i2\pi\nu_0 t}$. El espectro se mueve de 0-5 kHz a 100 MHz ± 5 kHz.

\textbf{¿Por qué queremos hacer esto?}
\begin{itemize}
\item Las antenas eficientes tienen tamaño relacionado con la longitud de onda
\item Para transmitir eficientemente por el aire, necesitamos frecuencias altas (longitudes de onda cortas)
\item Además, podemos tener múltiples usuarios en diferentes frecuencias portadoras sin que interfieran (multiplexación en frecuencia)
\end{itemize}

\textbf{En la práctica:}

Para señales reales usamos $\cos(2\pi\nu_0 t)$ o $\sin(2\pi\nu_0 t)$, que crean dos copias del espectro en $+\nu_0$ y $-\nu_0$ (frecuencias positivas y negativas).
\end{tcolorbox}

\subsection{Teorema de Parseval - Conservación de energía}

\begin{tcolorbox}[colback=verdeclaro!20,colframe=azuloscuro,title=\textbf{Teorema de Parseval}]
$$\boxed{E_s = \int_{\mathbb{R}} |s(t)|^2dt = \int_{\mathbb{R}} |S(\nu)|^2d\nu}$$

\textbf{Cada término significa:}
\begin{itemize}
\item $E_s$: energía total de la señal [Julios o voltios$^2$·segundo]
\item $|s(t)|^2$: potencia instantánea en el instante $t$
\item $|S(\nu)|^2$: densidad espectral de energía a la frecuencia $\nu$
\item Integral sobre $\mathbb{R}$: suma sobre todo el tiempo (izquierda) o toda la frecuencia (derecha)
\end{itemize}

\textbf{Significado profundo:}

La energía total de una señal es la misma si la calculamos en el dominio temporal o en el dominio frecuencial. La transformada de Fourier \textbf{preserva la energía}.

\textbf{Interpretación física:}

No importa si medimos la energía "segundo a segundo" en tiempo, o "frecuencia por frecuencia" en frecuencia. El total es el mismo.

\textbf{Utilidad práctica:}

Podemos calcular la energía de una señal en el dominio que sea más conveniente:
\begin{itemize}
\item Si $s(t)$ tiene forma simple → calcular en tiempo
\item Si $S(\nu)$ tiene forma simple → calcular en frecuencia
\end{itemize}

\textbf{Ejemplo:}

Para un pulso rectangular en tiempo, es más fácil calcular $\int |s(t)|^2 dt$ directamente. Para una señal que es suma de pocas frecuencias, es más fácil calcular $\int |S(\nu)|^2 d\nu$.
\end{tcolorbox}

\subsection{Identidades trigonométricas útiles}

Estas identidades son fundamentales para trabajar con señales moduladas (senos y cosenos).

\begin{tcolorbox}[colback=azulclaro!20,colframe=azuloscuro,title=\textbf{Identidades de potencias}]
$$\boxed{\cos^2(a) = \frac{1 + \cos(2a)}{2}}$$

$$\boxed{\sin^2(a) = \frac{1 - \cos(2a)}{2}}$$

\textbf{¿Para qué sirven?}

Cuando elevamos al cuadrado funciones seno o coseno (por ejemplo, al calcular potencia promedio), estas identidades permiten convertir $\cos^2$ o $\sin^2$ en una constante más un término con frecuencia doble.

\textbf{Ejemplo de aplicación:}

Si la señal es $s(t) = A\cos(2\pi f_c t)$, su potencia instantánea es:
$$P(t) = |s(t)|^2 = A^2\cos^2(2\pi f_c t) = A^2 \cdot \frac{1 + \cos(4\pi f_c t)}{2} = \frac{A^2}{2} + \frac{A^2}{2}\cos(4\pi f_c t)$$

Potencia promedio (promediando el término coseno que oscila):
$$P_{avg} = \frac{A^2}{2}$$
\end{tcolorbox}

\begin{tcolorbox}[colback=azulclaro!20,colframe=azuloscuro,title=\textbf{Productos de senos y cosenos}]
$$\boxed{\cos(a)\cos(b) = \frac{1}{2}[\cos(a-b) + \cos(a+b)]}$$

$$\boxed{\sin(a)\sin(b) = \frac{1}{2}[\cos(a-b) - \cos(a+b)]}$$

$$\boxed{\sin(a)\cos(b) = \frac{1}{2}[\sin(a+b) + \sin(a-b)]}$$

\textbf{¿Para qué sirven?}

Estas identidades son la base matemática de la \textbf{modulación y demodulación}.

\textbf{Ejemplo clave - Demodulación coherente:}

Señal recibida modulada: $r(t) = m(t)\cos(2\pi f_c t)$

Multiplicamos por la portadora local: $r(t) \cdot \cos(2\pi f_c t) = m(t)\cos^2(2\pi f_c t)$

Aplicando $\cos^2 = \frac{1+\cos(2\cdot)}{2}$:
$$= m(t) \cdot \frac{1 + \cos(4\pi f_c t)}{2} = \frac{m(t)}{2} + \frac{m(t)\cos(4\pi f_c t)}{2}$$

El primer término $\frac{m(t)}{2}$ es la señal deseada en banda base. El segundo término está a frecuencia $2f_c$ (muy alta) y se elimina con un filtro pasa-bajo.

\textbf{Resultado:} Recuperamos $m(t)$ (salvo un factor de escala).
\end{tcolorbox}

\begin{tcolorbox}[colback=azulclaro!20,colframe=azuloscuro,title=\textbf{Función sinc}]
$$\boxed{\text{sinc}(x) = \frac{\sin(\pi x)}{\pi x}}$$

\textbf{Propiedades importantes:}
\begin{itemize}
\item $\text{sinc}(0) = 1$ (por límite L'Hôpital)
\item $\text{sinc}(n) = 0$ para todo entero $n \neq 0$
\item Decae como $1/x$ para $|x|$ grande
\end{itemize}

\textbf{¿Dónde aparece?}

La función sinc es la transformada de Fourier del pulso rectangular:
$$\text{FT}[\text{rect}(t/T)] = T \cdot \text{sinc}(T\nu)$$

También aparece en el teorema de muestreo de Nyquist-Shannon: una señal muestreada se reconstruye como:
$$x(t) = \sum_{n} x[n] \cdot \text{sinc}(F_s t - n)$$

Es el pulso "ideal" de Nyquist: $\text{sinc}(t/T_s)$ vale 1 en $t=0$ y 0 en todos los múltiplos de $T_s$. Desafortunadamente, tiene duración infinita y no es realizable en práctica.
\end{tcolorbox}

\newpage

\section{Modulaciones Lineales}

\subsection{¿Por qué modular?}

\textbf{Problema:} Los bits son símbolos discretos (0 y 1), pero los canales físicos transmiten señales continuas (voltajes, campos electromagnéticos).

\textbf{Solución:} La \textbf{modulación} asigna a cada grupo de bits una forma de onda específica que se puede transmitir físicamente.

\textbf{Dos enfoques principales:}
\begin{enumerate}
\item \textbf{Banda base:} Transmitir directamente sin portadora (cables, Ethernet, USB)
\item \textbf{Paso de banda:} Modular sobre una portadora de alta frecuencia (radio, WiFi, 4G/5G, satélite)
\end{enumerate}

\subsection{Señal transmitida en banda base}

\begin{tcolorbox}[colback=verdeclaro!20,colframe=azuloscuro,title=\textbf{Señal banda base}]
$$\boxed{s(t) = \sum_{k=-\infty}^{\infty} a_k h(t-kT_s)}$$

\textbf{Cada símbolo significa:}
\begin{itemize}
\item $s(t)$: señal transmitida (función continua del tiempo)
\item $a_k$: símbolo de información en el instante $k$ (número real o complejo)
\item $h(t)$: pulso de forma o "shaping filter" (define cómo se "dibuja" cada símbolo)
\item $T_s$: período de símbolo [segundos] (tiempo entre símbolos consecutivos)
\item $k$: índice del símbolo (entero: ..., -1, 0, 1, 2, ...)
\item $h(t-kT_s)$: pulso $h(t)$ desplazado al instante $kT_s$
\item $\sum_{k=-\infty}^{\infty}$: suma sobre todos los símbolos
\end{itemize}

\textbf{Interpretación física:}

Cada símbolo $a_k$ "moldea" un pulso $h(t)$ que se transmite en su intervalo de tiempo correspondiente. La señal total es la superposición de todos estos pulsos desplazados.

\textbf{Ejemplo concreto:}

Si $a_0 = +1$, $a_1 = -1$, $a_2 = +1$, y $h(t)$ es un pulso rectangular de duración $T_s$:
\begin{itemize}
\item En $[0, T_s)$: transmitimos $+1 \cdot h(t)$ → voltaje positivo
\item En $[T_s, 2T_s)$: transmitimos $-1 \cdot h(t)$ → voltaje negativo
\item En $[2T_s, 3T_s)$: transmitimos $+1 \cdot h(t)$ → voltaje positivo
\end{itemize}

\textbf{¿Cuántos valores puede tomar $a_k$?}

Depende de la modulación:
\begin{itemize}
\item Binaria: $a_k \in \{-1, +1\}$ (2 valores → 1 bit por símbolo)
\item 4-PAM: $a_k \in \{-3, -1, +1, +3\}$ (4 valores → 2 bits por símbolo)
\item M-PAM: $M$ valores → $\log_2(M)$ bits por símbolo
\end{itemize}
\end{tcolorbox}

\subsection{Señal transmitida con portadora}

Cuando queremos transmitir por el aire (radio, WiFi, celular), necesitamos modular sobre una portadora de alta frecuencia.

\begin{tcolorbox}[colback=verdeclaro!20,colframe=azuloscuro,title=\textbf{Señal con portadora (forma real)}]
$$\boxed{s(t) = \sum_{k} a_k h(t-kT_s)\cos(2\pi\nu_0 t) - b_k h(t-kT_s)\sin(2\pi\nu_0 t)}$$

\textbf{Cada término significa:}
\begin{itemize}
\item $\nu_0$: frecuencia de la portadora [Hz]
\item $a_k$: componente "en fase" (I - In-phase) del símbolo $k$
\item $b_k$: componente "en cuadratura" (Q - Quadrature) del símbolo $k$
\item $\cos(2\pi\nu_0 t)$: portadora en fase
\item $-\sin(2\pi\nu_0 t)$: portadora en cuadratura (ortogonal a la anterior)
\item Las dos componentes se transmiten simultáneamente sin interferirse
\end{itemize}

\textbf{¿Por qué dos componentes?}

Las funciones $\cos(2\pi\nu_0 t)$ y $\sin(2\pi\nu_0 t)$ son \textbf{ortogonales}:
$$\int_0^{T} \cos(2\pi\nu_0 t)\sin(2\pi\nu_0 t)dt = 0$$

Esto significa que podemos transmitir información independiente en cada componente y recuperarlas por separado en recepción. Es como tener dos canales en uno.

\textbf{Ventaja clave:} Duplicamos la capacidad sin aumentar el ancho de banda.

\textbf{Ejemplo - QPSK:}
\begin{itemize}
\item $(a_k, b_k) = (+1, +1)$: transmite símbolo 00
\item $(a_k, b_k) = (+1, -1)$: transmite símbolo 01
\item $(a_k, b_k) = (-1, +1)$: transmite símbolo 10
\item $(a_k, b_k) = (-1, -1)$: transmite símbolo 11
\end{itemize}

Cada símbolo QPSK transmite 2 bits usando las dos dimensiones (I y Q).
\end{tcolorbox}

\begin{tcolorbox}[colback=amarillo!40,colframe=azuloscuro,title=\textbf{Notación compleja (equivalente)}]
$$\boxed{s_c(t) = \sum_{k} d_k h(t-kT_s)}$$

donde $\boxed{d_k = a_k + ib_k}$ (símbolo complejo)

La señal real se obtiene como:
$$\boxed{s(t) = \Re[s_c(t)e^{i(2\pi\nu_0 t + \theta_0)}]}$$

\textbf{Variables:}
\begin{itemize}
\item $s_c(t)$: envolvente compleja (señal banda base compleja)
\item $d_k = a_k + ib_k$: símbolo complejo (punto en el plano I-Q)
\item $i$: unidad imaginaria
\item $\theta_0$: fase inicial de la portadora
\item $\Re[\cdot]$: parte real
\end{itemize}

\textbf{¿Por qué usar notación compleja?}

Es mucho más compacta y elegante. En vez de trabajar con dos ecuaciones reales (componentes I y Q), trabajamos con una ecuación compleja. Toda la teoría de detección óptima, ecualización, etc., se simplifica enormemente.

\textbf{Relación con la forma real:}

Si desarrollamos $e^{i(2\pi\nu_0 t + \theta_0)} = \cos(2\pi\nu_0 t + \theta_0) + i\sin(2\pi\nu_0 t + \theta_0)$ y tomamos parte real, recuperamos la forma con seno y coseno.
\end{tcolorbox}

\textbf{Ventajas de usar portadora:}
\begin{enumerate}
\item \textbf{Adaptación al canal:} Transmitir en las frecuencias que se propagan mejor
\item \textbf{Tamaño de antenas:} Antenas eficientes tienen tamaño $\approx \lambda/4$ donde $\lambda = c/f$. Para frecuencias bajas, las antenas serían enormes.
\item \textbf{Multiplexación:} Varios usuarios en diferentes frecuencias portadoras sin interferirse
\item \textbf{Regulación:} Uso ordenado del espectro electromagnético
\end{enumerate}

\subsection{Parámetros fundamentales - Métricas de eficiencia}

\begin{tcolorbox}[colback=verdeclaro!20,colframe=azuloscuro,title=\textbf{Rapidez de modulación (Symbol Rate)}]
$$\boxed{D_s = \frac{1}{T_s} \quad \text{[bauds o símbolos/s]}}$$

\textbf{Variables:}
\begin{itemize}
\item $D_s$: rapidez de modulación (velocidad de símbolos)
\item $T_s$: período de símbolo [segundos]
\end{itemize}

\textbf{Significado:}

Es la velocidad a la que enviamos símbolos. Determina directamente el ancho de banda necesario.

\textbf{Ejemplo:} Si $T_s = 1$ ms, entonces $D_s = 1000$ símbolos/s = 1000 baudios.

\textbf{Nota importante:} No confundir con bits/s. Un símbolo puede representar múltiples bits.
\end{tcolorbox}

\begin{tcolorbox}[colback=verdeclaro!20,colframe=azuloscuro,title=\textbf{Débito binario (Bit Rate)}]
$$\boxed{D_b = D_s \cdot \log_2(M) = \frac{k}{T_s} \quad \text{[bits/s]}}$$

\textbf{Variables:}
\begin{itemize}
\item $D_b$: débito binario (velocidad de bits)
\item $M$: número de símbolos diferentes en el alfabeto
\item $k = \log_2(M)$: número de bits por símbolo
\end{itemize}

\textbf{Significado:}

Es la velocidad real de transmisión de información útil.

\textbf{Relación fundamental:} $D_b = D_s \cdot \log_2(M)$

\textbf{Ejemplos numéricos:}
\begin{itemize}
\item BPSK ($M=2$): $D_b = D_s \cdot 1 = D_s$ → 1 bit por símbolo
\item QPSK ($M=4$): $D_b = D_s \cdot 2 = 2D_s$ → 2 bits por símbolo
\item 16-QAM ($M=16$): $D_b = D_s \cdot 4 = 4D_s$ → 4 bits por símbolo
\item 64-QAM ($M=64$): $D_b = D_s \cdot 6 = 6D_s$ → 6 bits por símbolo
\end{itemize}

\textbf{Ejemplo concreto:}

Si transmitimos a $D_s = 1000$ símbolos/s usando 16-QAM:
$$D_b = 1000 \cdot \log_2(16) = 1000 \cdot 4 = 4000 \text{ bits/s} = 4 \text{ kbps}$$

\textbf{Formas de aumentar $D_b$:}
\begin{enumerate}
\item Aumentar $D_s$ (transmitir más rápido) → requiere más ancho de banda
\item Aumentar $M$ (más símbolos) → requiere más SNR (símbolos más cercanos)
\end{enumerate}
\end{tcolorbox}

\begin{tcolorbox}[colback=rojoclaro!20,colframe=azuloscuro,title=\textbf{Ancho de banda ocupado}]
\textbf{Para portadora única (QAM/PSK):}
$$\boxed{B_{occ} = \frac{1+\beta}{T_s} = (1+\beta)D_s}$$

\textbf{Para PAM/ASK:}
$$\boxed{B_{occ} = \frac{2(1+\beta)}{T_s} = 2(1+\beta)D_s}$$

\textbf{Variables:}
\begin{itemize}
\item $B_{occ}$: ancho de banda ocupado [Hz]
\item $\beta$: factor de roll-off del filtro (típicamente 0.25 a 0.5)
\item $(1+\beta)$: factor de exceso de banda
\end{itemize}

\textbf{Significado:}

Es el ancho de banda de frecuencias que necesitamos para transmitir nuestra señal.

\textbf{¿Por qué $(1+\beta)$?}

El ancho de banda mínimo teórico (criterio de Nyquist) es $B_{min} = \frac{1}{2T_s}$. En la práctica, usamos filtros raised cosine con factor de roll-off $\beta$, que ocupan:
$$B = \frac{1+\beta}{2T_s}$$

Para señales reales (AM), el espectro aparece en positivo y negativo → factor 2.

\textbf{Ejemplo:}

$D_s = 1000$ baudios, $\beta = 0.25$, QPSK:
$$B_{occ} = (1+0.25) \cdot 1000 = 1250 \text{ Hz}$$
\end{tcolorbox}

\begin{tcolorbox}[colback=amarillo!40,colframe=azuloscuro,title=\textbf{Eficiencia espectral}]
$$\boxed{\eta = \frac{D_b}{B_{occ}} \quad \text{[bits/s/Hz]}}$$

\textbf{Para portadora única:}
$$\boxed{\eta = \frac{\log_2(M)}{1+\beta}}$$

\textbf{Para PAM:}
$$\boxed{\eta = \frac{2\log_2(M)}{1+\beta}}$$

\textbf{Significado:}

Mide cuántos bits transmitimos por segundo por cada Hz de ancho de banda usado. Es la métrica clave de eficiencia espectral.

\textbf{Valores típicos en sistemas reales:}
\begin{itemize}
\item QPSK, $\beta=0.5$: $\eta = \frac{2}{1.5} \approx 1.33$ bits/s/Hz
\item 16-QAM, $\beta=0.25$: $\eta = \frac{4}{1.25} = 3.2$ bits/s/Hz
\item 64-QAM, $\beta=0.25$: $\eta = \frac{6}{1.25} = 4.8$ bits/s/Hz
\item LTE (4G): hasta 5-6 bits/s/Hz
\item 5G NR: hasta 7-8 bits/s/Hz
\end{itemize}

\textbf{Dilema fundamental:}
\begin{itemize}
\item Mayor $M$ → Mayor $\eta$ (más bits en mismo ancho de banda)
\item Pero mayor $M$ → Símbolos más cercanos → Mayor probabilidad de error
\item Solución: Necesitamos mayor SNR (más potencia de transmisión)
\end{itemize}

\textbf{Límite teórico - Fórmula de Shannon:}
$$\eta_{max} = \log_2\left(1 + \frac{S}{N}\right) \quad \text{[bits/s/Hz]}$$

Este es el máximo teórico alcanzable para un SNR dado. Ningún sistema real puede superarlo.
\end{tcolorbox}

Continuará con las secciones de alfabetos de modulación, Gray mapping, etc., todas explicadas en detalle...

\newpage

\section{Alfabetos de Modulación - Constelaciones}

[El documento continúa con explicaciones detalladas de cada constelación...]

\end{document}
